Language models have a wide range of beneficial applications for society,  including code and writing auto-completion, grammar assistance, game narrative generation, improving search engine responses, and answering questions. But they also have potentially harmful applications. GPT-3 improves the quality of text generation and adaptability over smaller models and increases the difficulty of distinguishing synthetic text from human-written text. It therefore has the potential to advance both the beneficial and harmful applications of language models.

Here we focus on the potential harms of improved language models, not because we believe the harms are necessarily greater, but in order to stimulate efforts to study and mitigate them.  The broader impacts of language models like this are numerous. We focus on two primary issues: the potential for deliberate misuse of language models like GPT-3 in Section \ref{section:misuse_of_language_models}, and issues of bias, fairness, and representation within models like GPT-3 in Section \ref{section:Fairness_Bias_and_Representation}.  We also briefly discuss issues of energy efficiency (Section \ref{section:Energy Usage}).
 
